\documentclass[10pt,a4paper]{amsart}
\usepackage[margin=19mm]{geometry}
\usepackage{amsmath,amssymb,mathtools}
\usepackage[scr=rsfs,cal=euler]{mathalfa}
\usepackage{xfrac}
\usepackage[unicode,hidelinks,bookmarks=false]{hyperref}
\newcommand{\ie}{i.e.\ }
\newcommand{\cf}{cf.\ }
\newcommand{\resp}{respectively\ }
\newcommand{\Subsection}{\subsection}
\providecommand{\nobreakdash}{\nobreak\hbox{-}}
\setlength{\emergencystretch}{2em}
\multlinegap0pt

\theoremstyle{plain}
\newtheorem{observation}{Observation}
\newtheorem*{obs}{Observation}
\theoremstyle{thmstyleone}
\newtheorem{theorem}{Theorem}
\newtheorem{proposition}[theorem]{Proposition}
\theoremstyle{thmstyletwo}
\newtheorem{example}{Example}
\newtheorem{remark}{Remark}
\newtheorem{lemma}{Lemma}
\newtheorem{corollary}{Corollary}
\theoremstyle{thmstylethree}
\newtheorem{definition}{Definition}
\title[Characterisations of Planar Galled Networks]{Characterisations of Planar Galled Networks}
\author{Hexuan Liu, Taoyang Wu, Guan-Ru Yu}
\date{Source: arXiv:2609.28597v2 (CC BY 4.0)}
\begin{document}
\maketitle

\noindent\emph{Source and licence.} The delivered work is the article shown in the title above, version arXiv:2609.28597v2 (CC BY 4.0), distributed under the Creative Commons Attribution 4.0 International licence, as recorded in the arXiv licence metadata for this exact version and shown on the abstract page saved with this item. The full licence notice, which carries the licence URL and the address of that abstract page, is the bundled file \texttt{licenses/arxiv-2609.28597-CC-BY-4.0.txt}.
\par\smallskip

\noindent\textit{Editorial scope.} Counted unit: the article's Theorem (source label noter2gall) (source0.tex lines 1118--1120, source label \texttt{noter2gall}): ``A galled network is terminal planar if and only if it contains no primitive tree vertex.''. It is delivered together with the article's complete original proof (source0.tex lines 1123--1174, one proof environment printed later in the article, detached from the statement by the article's own arrangement, 52 lines), copied line for line from the article's own text. Required context retained uncounted: lmm:cycle (lines 203--205); lmm:deg3-tree (lines 294--298); cycleandgalledcycle (lines 307--311); lmm:overlap-path-property (lines 788--797); lmm:cut-visible-basic (lines 913--921); lem:no\_s2 (lines 1033--1036); thm:normal-Kk23 (lines 1057--1059); lmm:cutvis-block (lines 1087--1093). Editorial formatting only: an AMS \texttt{amsart} preamble that restates the article's own macro definitions and its own theorem declarations, and the theorem environments the retained lines use; the article's own class file is neither shipped nor input; the retained environments are renumbered sequentially in order of appearance, whereas the article prints them with its own numbering; the article's equation numbering is not reproduced and every cross-reference inside the retained text resolves to the wrapper's numbering. No citation occurs inside any retained line, so the wrapper carries no bibliography; All retained bodies are exact copies of the bundled \texttt{sources/211602/source0.tex}, so no omission receipt applies. No asset-inclusion command occurs inside any retained span and the package ships no image asset.


\section*{Retained context: lmm:cycle (source0.tex lines 203--205)}

\begin{lemma}\label{lmm:cycle}
Let $N$ be a phylogenetic network. Then every cycle of $N$ contains at least one tree vertex and at least one reticulation vertex. Moreover, for any two tree vertices $u$ and $v$ of $N$, there is at most one path between $u$ and $v$ whose vertices are all tree vertices.
\end{lemma}


\section*{Retained context: lmm:deg3-tree (source0.tex lines 294--298)}

\begin{lemma}\label{lmm:deg3-tree}
Let $N$ be a galled network and let $B$ be a 2-connected subgraph of 
$\overline{N}$. Then every vertex 
$v \in V(B)$ with $\deg_B(v) = 3$ is a tree vertex of $N$.
\end{lemma}


\section*{Retained context: cycleandgalledcycle (source0.tex lines 307--311)}

\begin{theorem}\label{cycleandgalledcycle}
Let $C$ be a cycle in a galled network $N$ and suppose that $v$ is a vertex of $C$. Then there exists a reticulation vertex $r$ on $C$ such that $v$ is contained in the galled cycle $G_r$ associated with $r$. 
In particular, if $C$ contains exactly one reticulation vertex $r$, then 
$C=G_r$.
\end{theorem}


\section*{Retained context: lmm:overlap-path-property (source0.tex lines 788--797)}

\begin{lemma}\label{lmm:overlap-path-property}
Let $N$ be a galled network, and let $v$ be a primitive tree vertex that is an internal vertex of the intersecting path $P$ of two galled cycles $G_1$ and $G_2$. Let $H := G_1 \cup G_2$. Then we have the following three assertions:

(i) For any vertex $w \in V(H)$, every path in $N$ between $v$ and $w$ that is internally disjoint from $H$ has at least one reticulation vertex as an internal vertex.

(ii) $v$ is cut-visible from $H$ in $\overline{N}$.

(iii) There exists a path $P^*$ between $v$ and $x$ for $x\in \{\rho\}\cup X$ such that $P^*$ is internally disjoint from $H$.

\end{lemma}


\section*{Retained context: lmm:cut-visible-basic (source0.tex lines 913--921)}

\begin{lemma}\label{lmm:cut-visible-basic}
Let $H$ be a subgraph of a graph $G$, and let $B$ be a pivotal basis of $H$ in $G$. Then the following assertions hold:
\begin{enumerate}
\item[(i)] If $B'$ is a subset of $B$, then  $B'$ is also a pivotal basis of $H$ in $G$.
\item[(ii)] If $H'$ is a subgraph of $H$, then $V(H')\cap B$ is a pivotal basis of  $H'$ in $G$.
\item[(iii)] If a vertex $v$ in $H$ is a cut vertex of $G$,   then $B \cup \{v\}$ is also a pivotal basis of $H$ in $G$. 
\item[(iv)] If $G=\overline{N}$ for a galled network $N$, and $r$ is a reticulation vertex of $H$, then $B\cup\{r\}$ is also a pivotal basis of $H$ in $\overline{N}$.
\end{enumerate}
\end{lemma}


\section*{Retained context: lem:no\_s2 (source0.tex lines 1033--1036)}

\begin{lemma}\label{lem:no_s2}
A galled network $N$ is terminal planar if and only if $\overline{N}$ contains no pivotal subdivision of ${K}_{2:3}$.

\end{lemma}


\section*{Retained context: thm:normal-Kk23 (source0.tex lines 1057--1059)}

\begin{theorem}\label{thm:normal-Kk23}
Suppose $N$ is a galled network. Then $N$ is terminal planar if and only if $\overline{N}$ contains no normal pivotal subdivision of $K_{2:3}$.
\end{theorem}


\section*{Retained context: lmm:cutvis-block (source0.tex lines 1087--1093)}

\begin{lemma}\label{lmm:cutvis-block}
Let $H$ be a biconnected subgraph of a graph $G$ with $|V(H)|\ge 2$. 
Suppose $R=(v_0,v_1,\dots,v_k)$ is a pivotal ray of $H$ in $G$ with $k\ge 1$, and denote the block of $G$ containing $H$ by $B$.
Then $v_0$ is not a cut vertex of $G$, and the ray $R$ is contained in $B$ (i.e., $V(R)\subsetneq V(B)$). 
Furthermore, there exist a vertex $u\in V(H)\setminus\{v_0\}$ and a path $P$ between $v_0$ and $u$ such that $v_k\in V(P)$ and $V(P)\cap V(H)=\{v_0,u\}$.

\end{lemma}


\section{The counted unit: Theorem (source label noter2gall) and its complete original proof}

\begin{theorem}\label{noter2gall}
A galled network $N$ is terminal planar if and only if it contains no primitive tree vertex.
\end{theorem}

\begin{proof}
($\Rightarrow$) Suppose that $v$ is a primitive tree vertex of ${N}$, and let $G_1, G_2$ be two galled cycles such that $v$ is an internal vertex on their intersecting path $P$ with endpoints $t_1, t_2$. Let $r_1, r_2$ be the reticulations of $G_1, G_2$. Then $G_1 \cup G_2$ is a subdivision of $K_{2,3}$ in $\overline{N}$ with degree~$3$ branch vertices $t_1, t_2$, and degree~$2$ branch vertices $v, r_1, r_2$. By Lemma~\ref{lmm:overlap-path-property}(ii) and Lemma~\ref{lmm:cut-visible-basic}(iv), $\{v,r_1,r_2\}$ is a pivotal basis of $G_1\cup G_2$ in $\overline{N}$. Hence $G_1\cup G_2$ is a pivotal subdivision of $K_{2:3}$ in $\overline{N}$, and $N$ is not terminal planar by Lemma~\ref{lem:no_s2}.
\smallskip

($\Leftarrow$) Suppose $N$ is not terminal planar. By Theorem~\ref{thm:normal-Kk23}, $\overline{N}$ contains a normal pivotal subdivision $H$ of $K_{2:3}$ with branch vertices $t_1, t_2$ and three internally disjoint paths $P_1, P_2, P_3$. Without loss of generality, assume that $P_3$ is the unique tree path within $H$. Let $v$ be the branch vertex of $H$ corresponding to the vertex of $B$ that lies on $P_3$, and let $e=\{v,v'\}$ be the edge of $\overline{N}$ incident to $v$ such that $e \notin E(H)$. We show that ${N}$ contains a primitive tree vertex.

\smallskip
\noindent\textit{Case 1: $v$ is a cut vertex of $\overline{N}$.}

Applying Theorem~\ref{cycleandgalledcycle} to the cycles $P_1\cup P_3$ and $P_2\cup P_3$, we obtain reticulation vertices $r_1\in V(P_1)$ and $r_2\in V(P_2)$ whose galled cycles $G_1$ and $G_2$ both contain $v$. Since the cut edge $e$ is not in any cycle, $G_1$ and $G_2$ both contain the two incident edges of $v$ on $P_3$, so $v$ is an internal vertex on the intersecting path of $G_1$ and $G_2$, and hence $v$ is a primitive tree vertex.

\smallskip
\noindent\textit{Case 2: $v$ is not a cut vertex of $\overline{N}$.}

Note that $\overline{N}-v$ is connected. Let $\mathcal{P}$ be the set of paths in $\overline{N}$ between $v$ and some vertex of $V(H)\setminus\{v\}$ whose internal vertices do not lie in $H$. By Lemma~\ref{lmm:cutvis-block}, $\mathcal{P}$ is non-empty. Moreover, every path in $\mathcal{P}$ contains $e$, and both endpoints of every path in $\mathcal{P}$ are tree vertices by Lemma~\ref{lmm:deg3-tree}. We consider the following three cases.

\smallskip
\noindent\emph{Case 2(i): every path in $\mathcal{P}$ contains a reticulation vertex.}

By Theorem~\ref{cycleandgalledcycle} as in Case~1, there exist galled cycles $G_1$ and $G_2$, both containing $v$, with reticulation vertices on $P_1$ and $P_2$, respectively. If $e \in E(G_1)$, then $G_1$ contains a subpath $P$ in $\mathcal{P}$, which by assumption contains an internal reticulation vertex on $P$ and is distinct from $r_1$, contradicting that $G_1$ is a galled cycle. Similarly, $e\notin E(G_2)$. Hence both $G_1$ and $G_2$ contain the two incident edges of $v$ on $P_3$, and $v$ is a primitive tree vertex.



\smallskip
\noindent\emph{Case 2(ii): every path in $\mathcal{P}$ is a tree path.}

Let $R=(v_0,v_1,\dots,v_k)$ be the non-degenerate pivotal ray of $H$ in $\overline{N}$ with $v_0=v$ and $v_k=w$. By Lemma~\ref{lmm:cutvis-block}, there exists a path $P\in\mathcal{P}$ such that $w\in V(P)$. Note that $P$ is a tree path by assumption, hence $w$ is a tree vertex.  Let $u$ be the endpoint of $P$ in $V(H)\setminus\{v\}$.  If $u\in V(P_3)$, then $P \cup P_3[v,u]$ is a cycle consisting entirely of tree vertices, contradicting Lemma~\ref{lmm:cycle}. Without loss of generality, we may assume that $u\in V(P_1)$.

Consider the two cycles
$C_1=P\cup P_3[t_1,v]\cup P_1[t_1,u]$
and
$C_2=P\cup P_3[t_1,v]\cup P_2\cup P_1[u,t_2]$.
Both cycles contain $w$. Since $P$ and $P_3$ are tree paths, Lemma~\ref{lmm:cycle} and Theorem~\ref{cycleandgalledcycle} yield two galled cycles $G_1$ and $G_2$ containing $w$, whose associated reticulation vertices lie on $P_1[t_1,u]$ of $C_1$ and $P_1[u,t_2]\cup P_2$ of $C_2$, respectively.

Since $w$ is a cut vertex while $P \cup H$ is 2-connected, the edge incident with $w$ that does not belong to $P$ is a cut edge of $\overline{N}$, and hence cannot lie in either $G_1$ or $G_2$. Therefore $G_1$ and $G_2$ both contain the two incident edges of $w$ on $P$. Hence $w$ is an internal vertex on the intersecting path of $G_1$ and $G_2$, and so $w$ is a primitive tree vertex.





\smallskip
\noindent\emph{Case 2(iii): $\mathcal{P}$ contains both kinds of paths.}

Take $P^* \in \mathcal{P}$ containing a reticulation vertex. Let $r^*$ be the closest reticulation vertex to $v$ on $P^*$, and write
$P^*[v,r^*]=(w_0,\dots,w_k)$ with $w_0=v$ and $w_k=r^*$. Let $i$ be the largest index with $i<k$ such that $w_i$ lies on some tree path $P\in\mathcal{P}$. Note that $i\ne 0$; otherwise, $v'=w_1=r^*$ would be a reticulation vertex contained in every path in $\mathcal{P}$. Let $u$ be the endpoint other than $v$  of the tree path $P$. If $u\in V(P_3)$, then $P \cup P_3[u,v]$ is a cycle consisting entirely of tree vertices, contradicting Lemma~\ref{lmm:cycle}. Hence, without loss of generality, we may assume that $u\in V(P_1)$.

Consider the two cycles
$C_1=P_3[t_1,v]\cup P\cup P_1[t_1,u]$
and
$C_2=P_3[t_1,v]\cup P\cup P_1[u,t_2]\cup P_2$.
As in Case~2(ii), these cycles yield two galled cycles $G_1$ and $G_2$ containing $w_i$, whose associated reticulation vertices lie on $P_1[t_1,u]$ of $C_1$ and $P_1[u,t_2]\cup P_2$ of $C_2$, respectively. Note that $P_3[t_1,v]\cup P$ is a tree path with the internal tree vertex $w_i$. Set $H'=C_1\cup C_2$ and $e'=\{w_i,w_{i+1}\}$. By the maximality of $i$, every path from $w_i$ to $H'$ containing $e'$ contains a reticulation vertex. Hence the argument of Case~2(i) applies to $H'$ and $e'$, and $G_1$ and $G_2$ both contain the two incident edges of $w_i$ on $P$. Hence $w_i$ is a primitive tree vertex.
\end{proof}

\end{document}
